\def\topic{Analysis}
\ifdefined\MyHeaderIncluded
\else
\def\qwertyflag{true}

\ifdefined\MyHeaderIncluded

\else
  \newcommand{\foobar}{FooBar}
  \def\MyHeaderIncluded{true}
% if is till the end of the file

\documentclass[10pt,a4paper]{article}
\usepackage[utf8]{inputenc}
%\usepackage[german]{babel}
%\usepackage[cyrillic]{babel}
\usepackage[T1]{fontenc}
%\usepackage[T2A]{fontenc}
%\usepackage[english, russian]{babel}
\usepackage{amsmath}
\usepackage{amsfonts}
\usepackage{amssymb}
\usepackage{amsthm}
\usepackage{mathtools} %coloneqq

\usepackage{bussproofs} %

%\renewcommand{\thesubsection}{\arabic{subsection}}
%\newtheorem{thm}{Theorem}[section]
%\renewcommand{\thethm}{\arabic{subsection}.\arabic{thm}}

\usepackage{thmtools}
\usepackage{hyperref}
\hypersetup{
    colorlinks=true,
    linkcolor=blue,
    filecolor=magenta,      
    urlcolor=cyan,
    pdftitle={Overleaf Example},
    %pdfpagemode=FullScreen,
    }
%\usepackage{stix-type1}%\turnediota
\theoremstyle{definition}

\newtheorem{constructor}{Constructor}[section]
\newtheorem{iaxiom}{Axiom}
\newtheorem{axiom}{Axiom}
\newtheorem*{axiom*}{Axiom}
\newtheorem{definition}{Definition}[section]
\newtheorem*{definition*}{Definition}
\newtheorem{remark}{Remark}[section]
\newtheorem*{remark*}{Remark}
\newtheorem{recall}{Recall}[section]
\newtheorem*{recall*}{Recall}

\newtheorem{lemma}{Lemma}[section]
\newtheorem{theorem}{Theorem}[section]
\newtheorem*{theorem*}{Theorem}
\newtheorem{note}{Note}[section]
\usepackage{etoolbox}

%\newcommand*{\@Dashv}[2]{%
%  \reflectbox{$\m@th#1#2$}%
%}
\newcommand*{\separation}[3]{% 1)binder, 2)set, 3)condition
  \{#1\in #2: #3\}
}

% https://www.freecodecamp.org/news/simple-regex-tricks-for-beginners-3acb3fa257cb/
%  \\mathcal\{T\} |---> \topology
%  \\mathcal\{([^\}]*)\}   |--->  $1
\newcommand{\AND}{\ \ \&\ \ }
\newcommand{\any}{\mathrm{\ any}}
\newcommand{\borel}{\mathcal{B}}
\newcommand{\can}{\mathrm{can}}
\newcommand{\canonize}{\mathrm{canonize}}
\newcommand{\closure}{\mbox{cl}}
\newcommand{\complex}{\mathbb{C}}
\newcommand{\converges}{\longrightarrow}
\newcommand{\countable}{\mbox{\ countable}}
\newcommand{\ldefined}{\mbox{\ defined}}
\newcommand{\dom}{\mbox{dom}}
\newcommand{\Dom}{\mbox{Dom}}
\newcommand{\Done}{\noindent \textbf{Done.}\\}
\newcommand{\lexists}{\mbox{\ exists}}
\newcommand{\lforall}{\mbox{\ forall}}
\newcommand{\lunique}{\mbox{\ unique}}
\newcommand{\family}{\mathrm{\ family}}
\newcommand{\Fm}{\mathrm{Fm}}
\newcommand{\FuncSymb}{\mathrm{FuncSymb}}
\newcommand{\function}{\mathrm{\ function}}
\newcommand{\functional}{\mbox{\ functional}}
\newcommand{\IFF}{\ \Longleftrightarrow\ }
\renewcommand{\iff}{\ \longleftrightarrow\ }
\newcommand{\IMPLIES}{\ \Longrightarrow\ }
\renewcommand{\implies}{\ \longrightarrow\ }
\newcommand{\Inf}{\mathrm{Inf}}
\newcommand{\injective}{\mbox{\ injective}}
\newcommand{\injection}{\mbox{\ injection}}
\newcommand{\interior}{\mbox{int}}
\newcommand{\intersection}{\bigcap}
\newcommand{\isempty}{\mathrm{\ empty}}
\newcommand{\LB}{\mathrm{LB}}
\newcommand{\nghbset}{\mathcal{N}}
\newcommand{\nonempty}{\mathrm{\ nonempty}}
\newcommand{\Ord}{\mathrm{Ord}}
\newcommand{\power}{\mathcal{P}}
\newcommand{\pr}{\mathrm{pr}}
\newcommand{\fst}{\mathrm{fst}}
\newcommand{\snd}{\mathrm{snd}}
\newcommand{\PredSymb}{\mathrm{PredSymb}}
\newcommand{\Proof}{\noindent \textbf{Proof.}\\}
\newcommand{\ran}{\mathrm{ran}}
\newcommand{\Ran}{\mathrm{Ran}}
\newcommand{\real}{\mathbb{R}}
\newcommand{\relation}{\mathrm{\ relation}}
\newcommand{\set}{\mathrm{\ set}}
\newcommand{\singleton}{\mathrm{\ singleton}}
\newcommand{\semantics}{\mathrm{semantics}}
\newcommand{\semialgebra}{\mathcal{S}}
\newcommand{\Sup}{\mathrm{Sup}}
\newcommand{\surjective}{\mbox{\ surjective\ on\ }}
\newcommand{\Term}{\mathrm{Term}}
\newcommand{\Type}{\mathrm{Type}}
\newcommand{\topology}{\mathcal{T}}
\newcommand{\total}{\mbox{\ total\ on\ }}
\newcommand{\transitive}{\mathrm{\ transitive}}
\newcommand{\union}{\bigcup}
\newcommand{\lTrue}{\mathrm{True}}
\newcommand{\lFalse}{\mathrm{False}}
\newcommand{\lUncertain}{\mathrm{Uncertain}}
\newcommand{\lif}{\mathrm{if}}
\newcommand{\lthen}{\mathrm{then}}
\newcommand{\lelse}{\mathrm{else}}
\newcommand{\lexception}{\mathrm{exception}}
\newcommand{\lundef}{\mathrm{undef}}
\newcommand{\univocal}{\mathrm{\ univocal}}
\newcommand{\univocalat}{\mathrm{\ univocal\ at\ }}
\newcommand{\UB}{\mathrm{UB}}
\newcommand{\Un}{\mathcal{U}}
\newcommand{\Var}{\mathrm{Var}}

\newcommand{\notvarphi}{\phi}
\renewcommand{\phi}{\varphi}

\newcommand{\comment}[1]{\quad\mbox{(#1)}}
% local commands
%\preto\section{\undef\blipo} %etoolbox
\preto\subsection{\undef\blipo} %etoolbox
% https://tex.stackexchange.com/questions/105684/limit-macro-newcommand-scope-to-section
\usepackage{graphicx}
\usepackage{authblk}

\setlength\parindent{0pt} % whole doc of \noindent
\usepackage{soul}
%\usepackage{mathgl}

\fi % guard
\ifdefined\topic
\else
\def\topic{No Titel}
\fi
\title{\topic}
\author[1]{Georgy Dunaev}
\begin{document}
\maketitle
\fi
\section{Analysis}
This section develops topological limits along bases and their relation
to one-sided limits such as $f(x+0)$ and the familiar
$\varepsilon$--$\delta$ formulations. Amann and Escher's
\emph{Analysis I--III} provides the intended background context.

A further question is when convergence along a base can be tested by
sequences. Such a comparison needs explicit hypotheses on the spaces and
the base. The sequential test below supplies these hypotheses and gives
a counterexample to extending the result to arbitrary topological spaces.

\subsection{Limits along a proper filter base}
For the limit statements in this subsection, let $A$ be a set and
$\mathcal{B}$ a nonempty family of nonempty subsets of $A$, such that
for $B_1,B_2\in\mathcal{B}$ there is $B_3\in\mathcal{B}$ with
$B_3\subseteq B_1\cap B_2$. These are the local conditions for a
\emph{proper filter base}. They exclude the empty-family case admitted
by the exploratory filter axioms in the topology section.

Let $f:A\to Y$ be a function into a topological space. We write
$f\to y$ along $\mathcal{B}$ when
\[
\forall U\in\mathcal{N}_Y(y),\ \exists B\in\mathcal{B},\ f[B]\subseteq U.
\]
The graph of $f$ supplies a value at every point of $A$. If a formula
for $f$ originally denotes only a partial function, its domain must
first be restricted accordingly; a default value outside the domain
does not automatically belong in the limit calculation.

For a sequence $f:\mathbb{N}\to Y$, the tail sets
$\{n\in\mathbb{N}\mid n\geqslant N\}$ give ordinary sequential
convergence. For a right-hand limit at a real number $a$, use the
nonempty sets $A\cap(a,a+\delta)$, $\delta>0$, assuming that $a$ is
approached from the right by points of $A$. This specifies both the
side of approach and the domain on which evaluation is meaningful.

\begin{lemma}[Uniqueness of a limit in a Hausdorff space]
A function has at most one limit along a proper filter base when its
target space is Hausdorff.
\end{lemma}
\Proof
If distinct points $y,z$ were limits, choose disjoint neighbourhoods
$U$ of $y$ and $W$ of $z$. There would be base members $B,C$ with
$f[B]\subseteq U$ and $f[C]\subseteq W$. A nonempty base member
$D\subseteq B\cap C$ would then satisfy $f[D]\subseteq U\cap W
=\emptyclass$, which is impossible.
\Done

\subsection{Locality, composition, and real limits}
\begin{lemma}[Changing a function away from the approach region]
Let $f,g:A\to Y$, and suppose they agree on some $B_0$ in a proper
filter base $\mathcal B$. Then $f$ and $g$ have exactly the same limits
along $\mathcal B$.
\end{lemma}
\Proof
If $f\to y$ and $U$ is a neighbourhood of $y$, choose $B\in\mathcal B$
with $f[B]\subseteq U$. A base refinement $C\subseteq B\cap B_0$
satisfies $g[C]=f[C]\subseteq U$, so $g\to y$. Exchange $f,g$ for
the reverse implication.
\Done

\begin{proposition}[Pairing and continuous composition]
If $f:A\to Y$ and $g:A\to Z$ converge along the same proper filter
base to $y,z$, then $a\mapsto(f(a),g(a))$ converges to $(y,z)$ in
the product topology. If $h:Y\to Z$ is continuous at $y$ and $f\to y$,
then $h\circ f\to h(y)$.
\end{proposition}
\Proof
A neighbourhood of $(y,z)$ contains $U\times V$ for neighbourhoods
$U,V$ of $y,z$. Choose base members making $f$ lie in $U$ and $g$
lie in $V$, then a common refinement. On it the pair lies in the
product. For composition, continuity supplies a neighbourhood $U$ of
$y$ whose image lies in any prescribed neighbourhood of $h(y)$;
the limit of $f$ supplies a base member mapped into $U$.
\Done

\begin{theorem}[Punctured metric limits]
Let $A\subseteq X$, where $(X,d)$ is a metric space, let $a\in X$,
and assume every punctured ball about $a$ meets $A$. The sets
\[
B_\delta=\{x\in A:0<d(x,a)<\delta\},\qquad\delta>0,
\]
form a proper filter base. For a map $f:A\to(Y,e)$ into another metric
space and a proposed limit $y\in Y$, its limit along this base is $y$ exactly when
\[
\forall\varepsilon>0\ \exists\delta>0\ \forall x\in A,
\quad 0<d(x,a)<\delta\implies e(f(x),y)<\varepsilon.
\]
The value at $a$, when defined, has no effect on this limit.
\end{theorem}
\Proof
Nonemptiness is the approach assumption, and $B_{\min(\delta,\eta)}$
is contained in the intersection of two base members. Metric balls
form a neighbourhood base in $Y$, so the limit definition is exactly
the displayed quantifier condition. Since $a$ belongs to none of the
base members, changing its value does not change any image used here.
\Done

For a right-hand real limit, replace $0<d(x,a)<\delta$ by
$0<x-a<\delta$ and require that every such interval meets $A$.
The same proof gives the one-sided epsilon--delta criterion. Without
the approach assumption, an empty base member makes the quantified
condition vacuous for every proposed limit, which is why it cannot
be used to infer uniqueness at an isolated point.

\begin{theorem}[Squeezing along a proper base]
Let $f,h,g:A\to\mathbb R$ satisfy $f\leq h\leq g$ on some base
member. If both $f$ and $g$ tend to $L$ along the base, then $h\to L$.
\end{theorem}
\Proof
For $\varepsilon>0$, refine the ordering base member with base members
on which $f>L-\varepsilon$ and $g<L+\varepsilon$, respectively.
On the common refinement,
$L-\varepsilon<f\leq h\leq g<L+\varepsilon$, which is the limit
condition for $h$.
\Done

\subsection{When sequences test a limit}
A sequence $(B_n)$ of members of a filter base $\mathcal B$ is
\emph{cofinal} if every $B\in\mathcal B$ contains some $B_n$.
If such a sequence exists, put $C_n=\bigcap_{i=1}^n B_i$.
Finite induction using the refinement property shows that each $C_n$
contains a member of $\mathcal B$, and hence is nonempty. The decreasing
sequence $(C_n)$ is a new proper filter base, equivalent to the original:
every original member contains some $C_n$, and every $C_n$ contains an
original member. Consequently the two bases give exactly the same limits.
The sets $C_n$ need not themselves belong to the original base. This
construction requires no successive choices of refinements. The following
arguments use classical logic and countable choice for the explicitly
indexed nonempty sets of possible sequence terms.

\begin{theorem}[Sequential test for a countably based limit]
Let $(C_n)$ be decreasing and cofinal in a proper filter base on $A$.
For $f:A\to Y$ and $y\in Y$, the following are equivalent:
\begin{enumerate}
\item $f$ converges to $y$ along the base;
\item for every sequence $a_n\in C_n$, the sequence $f(a_n)$ converges
      to $y$ in $Y$.
\end{enumerate}
No countability or separation assumption on $Y$ is required.
\end{theorem}
\Proof
If $U$ is a neighbourhood of $y$, convergence gives a base member
$B$ with $f[B]\subseteq U$. Cofinality gives $C_N\subseteq B$;
decreasingness gives $a_n\in C_n\subseteq C_N$ for $n\geq N$.
Thus $f(a_n)$ is eventually in $U$.
Conversely, if the base limit fails, some neighbourhood $U$ of $y$
contains the image of no base member. Choose $a_n\in C_n$ with
$f(a_n)\notin U$ for every $n$. This sequence contradicts the second
condition.
\Done

\begin{corollary}[Sequential continuity at a first-countable point]
If $X$ has a countable neighbourhood base at $x$, then a function
$f:X\to Y$ is continuous at $x$ if and only if
$x_n\to x$ always implies $f(x_n)\to f(x)$.
\end{corollary}
\Proof
Continuity gives the forward assertion directly by inverse images of
neighbourhoods. For the converse, choose a decreasing neighbourhood
base $C_n$ at $x$. Every sequence $x_n\in C_n$ converges to $x$.
The theorem therefore gives convergence of $f$ along that base to
$f(x)$, which is precisely continuity at $x$.
\Done

\begin{example}[Why the countability hypothesis matters]
On an uncountable set $X$, give $X$ the cocountable topology: the
nonempty open sets are those with countable complement. A sequence
converging to $x$ is eventually equal to $x$. Indeed, the set of its
values different from $x$ is countable, so its complement is a
neighbourhood of $x$ that contains no other value of the sequence.
Consequently the identity from this topology to the discrete topology
on $X$ preserves all convergent sequences. It is not continuous,
because a singleton is discrete-open but not cocountable-open.
Thus preservation of sequences alone does not characterize continuity
on an arbitrary topological space.
\end{example}

\subsection{An elementary order argument}
\begin{theorem}[Nonnegativity from positive perturbations]
For a real number $a$,
\[
(\forall \varepsilon>0,\ a+\varepsilon\geqslant0)
\implies a\geqslant0.
\]
\end{theorem}
\Proof
Suppose $a<0$ and choose $\varepsilon=-a/2$. Then $\varepsilon>0$, but
\[
a+\varepsilon=a/2<0,
\]
contrary to the hypothesis. Hence $a\geqslant0$.
\Done
\ifdefined\qwertyflag
\undef\qwertyflag
\end{document}
\fi
